\chapter{Sonido, señal y percepción}\label{cap:sonido}

\enproceso{Versión condensada del capítulo 1 de la tesis (física del sonido, notas y frecuencias, timbre), más psicoacústica: percepción logarítmica de la altura, bandas críticas (Bark/ERB), sonoridad y enmascaramiento. Justifica por qué el eje de frecuencia de Fourierizer es logarítmico.}
